\section{Introduction}

La question de savoir si une baisse de la fiscalité du capital finance l'économie réelle ou
alimente simplement la distribution aux actionnaires structure l'évaluation des réformes fiscales
depuis deux décennies. La théorie des nouvelles vues de l'impôt sur les dividendes prédit que la
fiscalité de la distribution n'affecte pas l'investissement financé par les bénéfices retenus,
tandis que les anciennes vues prédisent un effet sur le coût du capital et donc sur
l'investissement \citep{auerbach1979,poterba1985,chetty2005}. L'arbitrage empirique entre ces
positions reste ouvert, et il dépend de ce que les entreprises et les ménages font effectivement de
l'argent qu'une réforme déplace \citep{hanlon2010}.

Le prélèvement forfaitaire unique, appelé flat tax dans le débat public, a été instauré par
l'article 28 de la loi de finances pour 2018, portée par le gouvernement issu de l'élection
présidentielle de 2017 \citep{loidefinances2018}. Il a remplacé l'imposition au barème progressif
de l'impôt sur le revenu par un taux proportionnel de 30 \%, soit 12,8 \% au titre de l'impôt sur le
revenu et 17,2 \% au titre des prélèvements sociaux, appliqué aux dividendes, aux intérêts et aux
plus-values de cession des personnes physiques résidentes. L'objectif affiché par le gouvernement
était double : simplifier une fiscalité jugée illisible et abaisser son taux marginal pour attirer
l'épargne vers le financement des entreprises et améliorer l'attractivité du pays. Le comité
d'évaluation des réformes de la fiscalité du capital, mandaté par les services du Premier ministre,
a suivi cette réforme et documenté la forte remontée des dividendes versés qui l'a accompagnée, sans
trancher sur l'effet réel \citep{francestrategie2023}. Ce choix d'une imposition proportionnelle du
capital n'est pas isolé en Europe. Il reprend la logique de l'impôt dual des pays nordiques, qui
taxent le revenu du capital à un taux proportionnel et le revenu du travail au barème progressif
depuis les réformes des années quatre-vingt-dix \citep{sorensen1994}, logique que l'Allemagne a
adoptée en 2009 avec son prélèvement libératoire, l'\textit{Abgeltungsteuer}, sur les revenus de
capitaux mobiliers. La réforme française met donc à l'épreuve, sur un cas net et récent, une
doctrine fiscale largement diffusée.

Ce papier pose la question de recherche suivante :

\begin{quote}
\itshape
Une baisse de la fiscalité des dividendes fait-elle entrer du capital neuf dans les entreprises, ou
se contente-t-elle de faire sortir des dividendes vers les ménages ?
\end{quote}

\noindent À cette question principale s'ajoute une question seconde, que la littérature a peu
traitée parce qu'elle exige de comparer deux chocs de sens opposé dans un même cadre : la réponse de
la distribution et de l'allocation du capital à l'impôt est-elle symétrique, de sorte qu'une baisse
défasse l'effet d'une hausse d'ampleur comparable, ou bien existe-t-il un cliquet, le capital
partant plus facilement qu'il ne revient ?

La France offre le terrain qui permet de répondre à ces deux questions, parce qu'elle a parcouru la
même marge fiscale dans les deux sens à cinq ans d'intervalle. Le prélèvement forfaitaire unique de
2018 a abaissé l'imposition des dividendes des personnes physiques résidentes, alors qu'en 2013 la
France avait suivi le chemin inverse en soumettant ces mêmes dividendes au barème progressif et en
alourdissant leur imposition. Ces deux réformes de sens opposé sur la même marge fournissent le
contraste dont la mesure d'un cliquet a besoin, là où les études antérieures n'observent qu'un choc
isolé.

La littérature nous arme pour lire ces chocs et pour situer notre apport. Sur le choc de 2013,
\citet{bach2023} montrent une contraction marquée des dividendes versés par les sociétés françaises,
concentrée sur les sociétés contrôlées par un actionnaire personne physique, et un déplacement de
l'arbitrage vers la rétention des bénéfices plutôt que vers l'investissement. Une première étude du
même choc avait rapporté une hausse de l'investissement \citep{boissel2022}, mais elle a été
rétractée après la mise en évidence de tendances pré-existantes entre sociétés traitées et de
contrôle \citep{bachcomment2023} ; l'ancrage robuste porte donc sur la rétention des bénéfices, non
sur l'investissement. Sur les baisses d'impôt sur les dividendes, les travaux de \citet{chetty2005}
sur la réforme américaine de 2003 documentent une hausse immédiate des distributions, tandis que
\citet{yagan2015} montre que la même baisse n'a relevé ni l'investissement ni l'emploi, un écart
entre la réponse de la distribution et celle de l'économie réelle qui annonce notre propre résultat.
\citet{alstadsater2017} établissent plus largement que la fiscalité de la distribution déplace
l'allocation du capital entre firmes sans nécessairement en changer le niveau agrégé, et
\citet{becker2013} montrent sur un panel international que les impôts sur les versements réorientent
l'investissement vers les firmes qui autofinancent leur croissance. Cette littérature converge vers
un même diagnostic : la fiscalité des dividendes agit d'abord sur le calendrier et la destination
des versements, non sur le niveau du capital productif. Notre contribution s'en distingue sur un
point. Là où chacun de ces travaux étudie un choc isolé, nous posons les deux chocs français dans un
cadre d'estimation unique et nous lisons l'écart de réponse, ce qui fait de l'asymétrie elle-même
l'objet de mesure, et non un sous-produit.

Nous apportons trois résultats. Premièrement, la baisse d'impôt de 2018 ne fait entrer aucun
capital neuf dans les fonds propres des sociétés, ni privées ni cotées, et ne soutient ni
l'investissement ni un recours accru à la dette. Deuxièmement, l'argent libéré sort vers les
ménages, dont les dividendes reçus bondissent de 32 \% l'année de la réforme, et se loge en épargne
liquide et en actions de portefeuille plutôt qu'en financement d'entreprise ou en immobilier.
Troisièmement, la comparaison des deux chocs révèle un cliquet, la baisse de 2018 ne défaisant qu'une
fraction de la réallocation de capital opérée par la hausse de 2013. Nous appuyons ces résultats par
un contrôle positif, la réplication de la hausse de la rétention des bénéfices postérieure à 2013
documentée par \citet{bach2023}, et par un contrôle négatif, l'absence de réponse là où le canal ne
peut pas jouer, les sociétés cotées à l'actionnariat dispersé.

La suite est organisée ainsi. La section~\ref{sec:cadre} développe le cadre théorique et les
hypothèses à partir d'un modèle d'arbitrage de l'actionnaire. La section~\ref{sec:donnees} décrit
les données et la construction des variables. La section~\ref{sec:modele} présente le modèle
empirique. La section~\ref{sec:resultats} rapporte les résultats, hypothèse par hypothèse. La
section~\ref{sec:robustesse} rassemble les tests de robustesse. La section~\ref{sec:conclusion}
conclut.
